\subsection{\texorpdfstring{CASE IN WHICH \(\mathbf{B}_{\mathrm{M}}\) IS POSITIVE AND DEGENERATE}{CASE IN WHICH B\_M IS POSITIVE AND DEGENERATE}}

In this number, we assume that S is finite, that \((W, S)\) is irreducible, and that the form \(B_{M}\) is positive and degenerate.

Lemma 2. The orthogonal complement \(\mathbf{E}^0\) of \(\mathbf{E}\) for \(\mathbf{B}_{\mathbf{M}}\) is of dimension 1; it is generated by an element \(v = \sum_{s\in S}v_s e_s\) with \(v_{s} > 0\) for all \(s\) .

This follows from Lemma 4 of §3, no. 5, applied to the matrix with entries \(\mathrm{B}_{\mathbb{M}}(e_s, e_{s'})\) .

Let \(v = \sum_{s} v_s e_s\) be the vector satisfying the conditions of Lemma 2 and such that \(\sum_{s} v_s = 1\) , and let A be the affine hyperplane of \(E^*\) consisting of the \(y^* \in E^*\) such that \(\langle v, y^* \rangle = 1\) . If T denotes the orthogonal complement of \(v\) in \(E^*\) , A has a natural structure of an affine space with space of translations T. Moreover, the form \(B_M\) defines by passage to the quotient a non-degenerate scalar product on \(E/E^0\) , hence also on its dual T; this gives a euclidean structure on the affine space A (Algebra, Chap. IX, §6, no. 6).

Let G be the subgroup of \(\mathbf{GL}(\mathbf{E})\) consisting of the automorphisms leaving v and \(B_{M}\) invariant; if \(g \in G\) , the contragredient map \({}^{t}g^{-1}\) leaves A and T stable, and defines by restriction to A a displacement \(i(g)\) of A (cf.~§3). It is immediate that this gives an isomorphism from G to the group of displacements of A. Moreover, the stabiliser \(G_{a}\) of a point a of A can be identified with the orthogonal group of the Hilbert space T and is thus compact. On the other hand, G is a locally compact group countable at infinity and A is a Baire space: it follows (Integration, Chap. VII, Appendix, Lemma 2) that the map \(\psi : g \mapsto g(a)\) defines a homeomorphism from \(G/G_a\) to A. Thus G acts properly on A (General Topology, Chap. III, §4, no. 2, Prop. 5). Since W is a subgroup of G, it can be identified with a group of displacements of A. We are going to show that this group satisfies the assumptions of §3. More precisely:

PROPOSITION 10. The group W provided with the discrete topology acts properly on A; it is generated by orthogonal reflections; it is infinite, irreducible and essential (§3, no. 7). The intersection C ∩ A is a chamber of A for W. If L\_s denotes the hyperplane of A formed by the intersection of A with the hyperplane of E* orthogonal to e\_s, the L\_s for s ∈ S are the walls of C ∩ A. If ε\_s is the unit vector of T orthogonal to L\_s on the same side of L\_s as C ∩ A, then (ε\_s\textbar ε\_t) = -cos(π/m(s,t)) (for s,t ∈ S) and the Coxeter matrix of W (§3, no. 4) is identified with M.

By Cor. 3 of Th. 1, W is discrete in \(\mathbf{GL}(\mathbf{E})\) , and hence in \(\mathbf{G}\) , and acts properly on A. Let \(s \in S\) . Since \(\operatorname{Card} S \geqslant 2\) , the hyperplane of \(\mathbf{E}^*\) orthogonal to \(e_s\) is not orthogonal to \(v\) and its intersection \(\mathbf{L}_s\) with A is indeed a hyperplane. The displacement corresponding to \(s\) is thus a displacement of order 2 leaving fixed all the points of \(\mathbf{L}_s\) : it is necessarily the orthogonal reflection associated to \(\mathbf{L}_s\) . It follows that W is generated by orthogonal reflections. Th. 2 now shows that it is infinite and Prop. 7 that it is essential and irreducible.

Since C is an open simplicial cone, whose walls are the hyperplanes with equations \(\langle x^{*}, e_{s} \rangle = 0\) (for \(s \in S\) ), the intersection \(C \cap A\) is a convex, hence connected, open and closed subset of the complement of the union of the \(L_{s}\) in A. Moreover, \(C \cap A\) is non-empty, for if \(x^{*} \in C\) we have \(\langle x^{*}, v \rangle = \sum_{s} v_{s} \langle x^{*}, e_{s} \rangle > 0\) and \(\langle x^{*}, v \rangle^{-1} x^{*} \in C \cap A\) . It follows that \(C \cap A\) is a chamber of A relative to the system of the \(L_{s}\) . Moreover, \(w(C \cap A) \cap L_{s} = \varnothing\) for all \(w \in W\) (cf.~no. 4, property \((\mathrm{P}_{n})\) ) and it follows that \(C \cap A\) is a chamber of A relative to the system consisting of the transforms of the \(L_{s}\) by the elements of W; by the Cor. of Th. 1 of §3, no. 2, it follows that \(C \cap A\) is a chamber of A relative to W.

Let \(a_{s}^{*}\) be the vertex of the simplex \(C \cap A\) not in \(L_{s}\) . We have

\[
\left\langle a _ {s} ^ {*}, e _ {t} \right\rangle = 0
\]

for \(s,t\in \mathrm{S}\) and \(s\neq t\) , and

\[
\langle a _ {s} ^ {*}, e _ {s} \rangle = v _ {s} ^ {- 1} \langle a _ {s} ^ {*}, v \rangle = v _ {s} ^ {- 1}.
\]

Let \(\varepsilon_{s}\) be the vector in T defined by the relations:

\[
\left(\varepsilon_ {s} \mid a _ {s} ^ {*} - a _ {t} ^ {*}\right) = v _ {s} ^ {- 1} \text {for} t \in S, t \neq s.
\]

The vector \(\varepsilon_{s}\) is orthogonal to \(L_{s}\) and is on the same side of the hyperplane \(L_{s}\) as \(C \cap A\) . Moreover

\[
\left(\varepsilon_ {s} \mid a _ {s} ^ {*} - a _ {t} ^ {*}\right) = \left\langle e _ {s}, a _ {s} ^ {*} - a _ {t} ^ {*} \right\rangle \text {for all} s, t \in S
\]

which shows that \(\varepsilon_{s}\) is the image of the class of \(e_{s}\) under the isomorphism from \(E/E^{0}\) to T given by the quadratic form \(B_{M}\) . It follows that

\[
(\varepsilon_ {s} | \varepsilon_ {t}) = \mathrm{B} _ {\mathrm{M}} (e _ {s}, e _ {t}).
\]

Consequently \(\varepsilon_{s}\) is a unit vector and the last assertion of Prop. 10 is proved.

The euclidean affine space A provided with the group W is called the space associated to the Coxeter matrix M and we denote it by \(A_{M}\) .

Proposition 10 admits a converse:

PROPOSITION 11. Let W be a group of displacements of a euclidean affine space A, satisfying the assumptions of §3. Assume that W is infinite, essential and irreducible. Then the form \(B_{M}\) attached to the Coxeter matrix M of W is positive degenerate and there exists a unique isomorphism from the affine space \(A_{M}\) associated to M to A, commuting with the action of W. This isomorphism transforms the scalar product of \(A_{M}\) into a multiple of the scalar product of A.

Let \(C_0\) be a chamber of A and let S be the set of orthogonal reflections with respect to the walls of \(C_0\) . If \(\eta_s\) denotes the unit vector orthogonal to the hyperplane \(N_s\) associated to the reflection s and on the same side of \(N_s\) as \(C_0\) (§3, Prop. 3), the form \(B_M\) is such that \(B_M(e_s, e_t) = (\eta_s | \eta_t)\) for \(s, t \in S\) . It is thus positive. Since the \(\eta_s\) are linearly dependent (§3, no. 9, Prop. 8), it is degenerate.

We can thus apply the preceding constructions to M. With the same notation as above, \((\varepsilon_{s}|\varepsilon_{t}) = (\eta_{s}|\eta_{t})\) and there exists a unique isomorphism \(\varphi\) of Hilbert spaces from T to the space of translations of A such that \(\varphi(\varepsilon_{s}) = \eta_{s}\) . Let a and b be two distinct vertices of \(C_{0}\) and \(s_{0}\) the reflection in S such that \(a \notin N_{s_{0}}\) . Put \(\lambda = (\eta_{s}|a - b)\) and let \(\psi\) be the affine bijection from \(A_{M}\) to A defined by

\[
\psi \left(a _ {s _ {0}} + x\right) = a + v _ {s _ {0}} \lambda \varphi (x) \text {for} x \in \mathrm{T}.
\]

It is then immediate that \(\psi(\mathrm{L}_{s}) = \mathrm{N}_{s}\) for all \(s \in S\) and that \(\psi\) transforms the scalar product of \(A_{M}\) into a multiple of that on A. It follows at once that \(\psi\) commutes with the action of W. Finally, the uniqueness of \(\psi\) is evident, for \(a_{s}\) for example is the unique point of \(A_{M}\) invariant under the reflections \(t \in S, t \neq s\) .

